% Supplementary Material — standalone \input{} file
% Requires: amsmath, amssymb, booktabs, array, float, enumitem packages in parent

\section{Extended Derivations}
\label{app:derivations}

\subsection{Geodesic distance on the Grassmannian}
\label{app:geodesic}

The Grassmannian $\Gr(k,d)$ admits a canonical Riemannian metric inherited
from the ambient space of $d \times k$ orthonormal frames. Given two
subspaces $[Q_1], [Q_2] \in \Gr(k,d)$ represented by orthonormal bases
$Q_1, Q_2 \in \R^{d \times k}$, the principal angles
$0 \leq \theta_1 \leq \cdots \leq \theta_k \leq \pi/2$ are defined via
the singular value decomposition
\begin{equation}
Q_1^\top Q_2 = U \, \diag(\cos\theta_1, \ldots, \cos\theta_k) \, V^\top.
\label{eq:principal_angles_svd}
\end{equation}
The geodesic distance under the canonical metric is
\begin{equation}
\dgr([Q_1],[Q_2]) = \left(\sum_{i=1}^k \theta_i^2\right)^{1/2}.
\label{eq:geodesic_full}
\end{equation}
This distance is gauge-invariant: for any $R_1, R_2 \in O(k)$,
$\dgr([Q_1 R_1], [Q_2 R_2]) = \dgr([Q_1], [Q_2])$, because replacing
$Q_j$ by $Q_j R_j$ yields the same column span. The maximum distance
between two $k$-dimensional subspaces is $\pi\sqrt{k}/2$, achieved when
all principal angles equal $\pi/2$ (orthogonal subspaces).

For random subspaces drawn uniformly from the Haar measure on $\Gr(k,d)$,
the expected pairwise distance concentrates around
$\pi\sqrt{k}/2 \cdot \sqrt{1 - k/d}$ when $k \ll d$
\citep{edelman1998geometry, absil2008optimization}. In our setting
($k = 2$, $d = 128$), this gives an expected random distance of
$\approx 2.20$ radians, consistent with the observed inter-seed distances
of $[1.93, 2.19]$ (\S\ref{sec:gauge_freedom}).

\subsection{Equivariance test statistic}
\label{app:equiv_test}

For an operation $f(a,b) = y$ with a cyclic group action
$a \mapsto a + g \pmod{p}$, we test whether the DAS subspace respects
the group symmetry via an \emph{intervention-based} test. Let
$Q \in \R^{d \times k}$ be the learned DAS rotation. The equivariance
condition requires that shifting input $a$ by $g$ shifts the model's
output by $g$ (modulo the operation's output range $p_\text{mod}$):
\begin{equation}
f(a + g, b) \equiv f(a, b) + g \pmod{p_\text{mod}}
\label{eq:equiv_condition}
\end{equation}
where $p_\text{mod} = p$ for pure modular operations and
$p_\text{mod} = \min(91, p{-}1)$ for composite addition
($a + b \bmod 91$ on a field of size $p$).

We test this by performing DAS-style interventions. For each valid
triple $(a, b, g)$:
\begin{enumerate}[leftmargin=*]
\item Compute the base activation $h(a, b)$ and the shifted activation
  $h(a{+}g, b)$.
\item Replace the in-subspace component: $h' = h(a,b) + Q Q^\top
  \bigl(h(a{+}g, b) - h(a, b)\bigr)$.
\item Run the model's remaining layers on $h'$ and obtain the predicted
  label $\hat{y}'$.
\item Count as equivariant if $\hat{y}' = (y + g) \bmod p_\text{mod}$,
  where $y = f(a, b)$.
\end{enumerate}
The reported equivariance fraction is
\begin{equation}
\text{Equiv} = \frac{1}{|\mathcal{T}|} \sum_{(a,b,g) \in \mathcal{T}}
\mathbf{1}\bigl[\hat{y}' = (y + g) \bmod p_\text{mod}\bigr]
\label{eq:equiv_fraction}
\end{equation}
where $\mathcal{T}$ is the set of all valid (input, shift) triples on the
test set. This test directly measures whether the learned subspace
captures the group action's causal structure: swapping the
in-subspace component between shifted inputs should produce the
correspondingly shifted output. Random orthogonal subspaces achieve
equivariance $< 1\%$ across all operations tested.

\paragraph{Relationship to angular tolerance test.}
An earlier version of this metric estimated the rotation angle
$\hat\phi_g$ induced by each shift $g$ in the projected space and
counted an input as equivariant if its angular displacement fell
within $\varepsilon = 2\pi/p$ of $\hat\phi_g$. This angular tolerance
becomes prohibitively tight for large primes ($\varepsilon = 0.003$
radians at $p = 211$), producing artificially low equivariance.
The intervention-based test avoids this issue by directly testing the
causal consequence of the group action.

\subsection{Diversity ratio}
\label{app:diversity_ratio}

The intervention diversity ratio $\rho$ (Eq.~\ref{eq:diversity_ratio})
measures whether an encoder-decoder intervention preserves input-specific
information or collapses to a class-conditional lookup table.

For a set of intervened activations $\{h'_i\}$ sharing source label $y_s$,
the within-class standard deviation is
$\sigma_{y_s} = \bigl(\frac{1}{n_{y_s}} \sum_i \|h'_i - \bar{h}'_{y_s}\|^2
\bigr)^{1/2}$
where $\bar{h}'_{y_s}$ is the mean intervened activation for source label $y_s$.
The denominator uses the same statistic computed on natural (non-intervened)
activations sharing the same label. The ratio $\rho = \mathbb{E}_{y_s}[\sigma_{y_s}^{\text{iv}}]
/ \mathbb{E}_{y_s}[\sigma_{y_s}^{\text{nat}}]$ ranges from 0 (complete collapse: all
interventions for the same source produce the same activation) to
$\approx 1$ (intervention preserves the natural within-class variability).

Values of $\rho \gg 1$ are possible in principle (the intervention adds noise),
but in practice indicate that the decoder is amplifying base-specific variation.
The observed NL-DAS values ($\rho \approx 0.05$ on IOI) confirm near-complete
collapse; the Structured pi-SAE values ($\rho \approx 0.89$) confirm
preservation of nuisance information.


\section{Hyperparameter Tables}
\label{app:hyperparameters}

\subsection{Grokking model training}

\begin{table}[H]
\centering
\small
\setlength{\tabcolsep}{4pt}
\begin{tabular}{ll}
\toprule
\textbf{Parameter} & \textbf{Value} \\
\midrule
Architecture & 1-layer transformer \\
$d_\text{model}$ & 128 \\
$n_\text{heads}$ & 4 \\
$d_\text{head}$ & 32 \\
$d_\text{MLP}$ & 512 \\
Optimizer & AdamW \\
Learning rate & $10^{-3}$ \\
Weight decay & 1.0 \\
Batch size & Full batch \\
Train/test split & 30\% / 70\% \\
Training epochs & 25{,}000--80{,}000 (operation-dependent) \\
Deeper models & 1.5$\times$ epochs for 2L and 4L \\
Grokking criterion & Test cross-entropy $< 0.1$ \\
Moduli tested & $p \in \{53, 97, 113, 211\}$ \\
Depths tested & $n_\text{layers} \in \{1, 2, 4\}$ \\
\bottomrule
\end{tabular}
\caption{Grokking model training hyperparameters. All operations use the same
architecture and optimizer; only epoch count varies.}
\label{tab:hp_grokking}
\end{table}

\subsection{DAS fitting}

\begin{table}[H]
\centering
\small
\begin{tabular}{ll}
\toprule
\textbf{Parameter} & \textbf{Value} \\
\midrule
Optimizer & Adam \\
Learning rate & $10^{-3}$ \\
Gradient steps & 200 \\
Batch size & 16 interchange pairs \\
$k$ values swept & $\{2, 4, 6, 8, 10, 12, 16, 20, 24, 32\}$ \\
Initialization & Random orthonormal $Q \in \R^{d \times k}$ \\
Intervention site & Residual stream (post-MLP, layer 0) \\
Hard-example selection & $y_b \neq y_s$, both correctly classified \\
\bottomrule
\end{tabular}
\caption{DAS fitting hyperparameters. Each $k$ value is an independent
optimization, not a nested subspace.}
\label{tab:hp_das}
\end{table}

\subsection{Structured pi-SAE}

\begin{table}[H]
\centering
\small
\begin{tabular}{ll}
\toprule
\textbf{Parameter} & \textbf{Value} \\
\midrule
Encoder & 2-layer MLP ($d_\text{model} \to 128 \to z$), ReLU \\
Decoder & 2-layer MLP ($z \to 128 \to d_\text{model}$), ReLU \\
$k$ (causal latent dim) & $\{2, 4, 8\}$ \\
$m$ (nuisance latent dim) & 15 \\
Sparsity & $L_1$ on $z_\text{causal}$, expansion $8\times$ \\
Prior & $p(z_\text{causal} \mid y) = \mathcal{N}(\mu_y, \sigma_y^2 I)$ \\
$\beta$ (KL weight) & 1.0 \\
$\alpha$ (classification weight) & 10.0 \\
$\lambda$ ($L_1$ coefficient) & 1.0 \\
Optimizer & Adam, lr $= 10^{-3}$ \\
Training epochs & 300 (reconstruction + classification) \\
E2E phase & 200 warmup + E2E intervention CE \\
E2E intervention & Additive: $h' = h_b + \text{dec}(z_\text{swap}) - \text{dec}(z_\text{orig})$ \\
\bottomrule
\end{tabular}
\caption{Structured pi-SAE hyperparameters. The E2E phase is used for
GPT-2 language tasks; grokking models use reconstruction + classification only.}
\label{tab:hp_pisae}
\end{table}

\subsection{NL-DAS baseline}

\begin{table}[H]
\centering
\small
\begin{tabular}{ll}
\toprule
\textbf{Parameter} & \textbf{Value} \\
\midrule
Encoder $f$ & 2-layer MLP ($d \to 256 \to d$), ReLU \\
Decoder $g$ & 2-layer MLP ($d \to 256 \to d$), ReLU \\
$Q$ & Learned jointly, $\R^{d \times k}$ \\
Training objective & Interchange loss only (no reconstruction) \\
NL-DAS+recon variant & $+ \lambda_\text{recon} \|g(f(h)) - h\|^2$, $\lambda = 0.1$ \\
Optimizer & Adam, lr $= 10^{-3}$ \\
Training steps & 200 \\
\bottomrule
\end{tabular}
\caption{NL-DAS hyperparameters. Architecture matches the pi-SAE encoder
to ensure a fair capacity comparison.}
\label{tab:hp_nldas}
\end{table}


\section{Per-Operation Atlas Data}
\label{app:atlas_data}

Table~\ref{tab:full_atlas} reports the complete atlas measurements for all
14 operations at $p = 113$ (1-layer), including metrics omitted from the
main text for space.

\vspace{1em}
\begin{table}[H]
\centering
\small
\setlength{\tabcolsep}{3pt}
\renewcommand{\arraystretch}{1.15}
\begin{tabular}{@{}lcccccccc@{}}
\toprule
\textbf{Operation} & \textbf{Grok} & \textbf{Loss} & \textbf{IIA}$_{k{=}2}$ &
\textbf{IIA}$_{k{=}4}$ & \textbf{IIA}$_{k{=}8}$ & $k^*$ &
\textbf{Equiv.} & \textbf{Class} \\
\midrule
\multicolumn{9}{l}{\textit{Always Grassmannian}} \\
Multiplication   & \checkmark & 0.000 & 1.000 & 1.000 & 1.000 & 2 & 0.995 & Always \\
Subtraction      & \checkmark & 0.000 & 1.000 & 1.000 & 1.000 & 2 & 1.000 & Always \\
Division         & \checkmark & 0.000 & 1.000 & 1.000 & 1.000 & 2 & 1.000 & Always \\
Bitwise XOR      & \checkmark & 0.003 & 1.000 & 1.000 & 1.000 & 2 & 1.000 & Always \\
Cubic sum        & \checkmark & 0.000 & 1.000 & 1.000 & 1.000 & 2 & 1.000 & Always \\
Sum of squares   & \checkmark & 0.005 & 1.000 & 1.000 & 1.000 & 2 & 0.987 & Always \\
Max              & \checkmark & 0.074 & 0.960 & 1.000 & 1.000 & 2 & 0.957 & Always \\
\midrule
\multicolumn{9}{l}{\textit{Stochastic}} \\
Comp.\ add.\ (orig) & \checkmark & 0.000 & --- & --- & --- & --- & 0.998 & Stochastic \\
Comp.\ add.\ (s42)  & $\times$   & 10.21 & 0.800 & 0.860 & 0.900 & 6 & 0.110 & Stochastic \\
Power (s137)     & \checkmark & 0.088 & 0.940 & 1.000 & 1.000 & 4 & 0.965 & Stochastic \\
Power (s42)      & $\times$   & 1.143 & 0.980 & 1.000 & 1.000 & 2 & 0.665 & Stochastic \\
\midrule
\multicolumn{9}{l}{\textit{Never Grassmannian}} \\
Cubing           & $\times$ & 9.765 & 0.857 & 1.000 & 1.000 & 4 & 0.509 & Never \\
Squaring         & $\times$ & 7.808 & 0.857 & 1.000 & 1.000 & 4 & 0.314 & Never \\
Abs.\ difference & $\times$ & 1.835 & 0.800 & 0.840 & 0.860 & $>$32 & 0.189 & Never \\
Polynomial       & $\times$ & 21.55 & 0.720 & 0.780 & 0.820 & $>$32 & 0.015 & Never \\
Affine           & $\times$ & 26.21 & 0.780 & 0.800 & 0.840 & $>$32 & 0.026 & Never \\
\bottomrule
\end{tabular}
\renewcommand{\arraystretch}{1.0}
\caption{Complete atlas at $p = 113$ (1-layer). IIA is reported at $k \in \{2, 4, 8\}$;
$k^*$ is the smallest $k$ achieving IIA $\geq 0.95$. Grassmannian operations
saturate by $k = 2$; non-Grassmannian operations show gradual growth.
Equivariance is the fraction of test inputs satisfying the rotation
property (\S\ref{app:equiv_test}).}
\label{tab:full_atlas}
\end{table}

Table~\ref{tab:seed_data} reports per-seed DAS results for the 10 seeds
of modular multiplication used in the subspace degeneracy analysis
(\S\ref{sec:gauge_freedom}).

\vspace{1em}
\begin{table}[H]
\centering
\small
\begin{tabular}{ccccc}
\toprule
\textbf{Seed} & \textbf{Test loss} & \textbf{IIA}$_{k{=}2}$ &
\textbf{Equiv.} & \textbf{Mean $d_\Gr$ to others} \\
\midrule
42   & 0.000 & 1.000 & 0.995 & 2.08 \\
137  & 0.000 & 0.980 & 0.991 & 2.11 \\
256  & 0.000 & 1.000 & 0.998 & 2.05 \\
512  & 0.000 & 0.960 & 0.989 & 2.09 \\
1024 & 0.000 & 1.000 & 0.997 & 2.07 \\
2024 & 0.000 & 0.940 & 0.985 & 2.13 \\
3000 & 0.000 & 1.000 & 0.993 & 2.06 \\
4096 & 0.000 & 0.980 & 0.990 & 2.10 \\
5555 & 0.000 & 1.000 & 0.996 & 2.08 \\
9999 & 0.000 & 1.000 & 0.992 & 2.09 \\
\midrule
\multicolumn{2}{l}{Mean $\pm$ std} & $0.986 \pm 0.020$ & $0.993 \pm 0.004$ &
$2.09 \pm 0.02$ \\
\bottomrule
\end{tabular}
\caption{Per-seed results for modular multiplication ($p = 113$, 1-layer).
All 10 seeds grok and achieve high IIA and equivariance, but their DAS
subspaces are near-orthogonal (mean pairwise $d_\Gr = 2.09$, compared to
the random-subspace expectation of $\approx 2.20$ on $\Gr(2, 128)$).}
\label{tab:seed_data}
\end{table}


\section{Pi-SAE Ablation Details}
\label{app:ablation}

The $2 \times 2$ ablation (Table~\ref{tab:ablation}) tests whether both
components of the Structured pi-SAE---the label-conditional prior and the
sparsity constraint---are individually necessary. This appendix provides
the full per-operation breakdown including reconstruction MSE, which
serves as a built-in falsifiability diagnostic.

\vspace{1em}
\begin{table}[H]
\centering
\small
\setlength{\tabcolsep}{3.5pt}
\renewcommand{\arraystretch}{1.15}
\begin{tabular}{@{}l|cc|cc|cc|cc@{}}
\toprule
& \multicolumn{2}{c|}{\textbf{VAE}} & \multicolumn{2}{c|}{\textbf{SAE}}
& \multicolumn{2}{c|}{\textbf{pi-VAE}} & \multicolumn{2}{c}{\textbf{Str.\ pi-SAE}} \\
\textbf{Operation} & IIA & MSE & IIA & MSE & IIA & MSE & IIA & MSE \\
\midrule
\multicolumn{9}{l}{\textit{Grokked operations}} \\
Addition        & 0.02 & 0.8 & 1.00 & 1.1  & 0.00 & 0.7  & \textbf{1.00} & 0.6 \\
Multiplication  & 0.00 & 0.9 & 0.72 & 1.3  & 0.00 & 0.8  & \textbf{1.00} & 0.7 \\
Quartic sum     & 0.02 & 1.0 & 0.32 & 1.5  & 0.00 & 0.9  & \textbf{1.00} & 0.8 \\
IOI (GPT-2)     & 0.56 & 2.1 & 0.83 & 3.4  & 0.95 & 1.8  & \textbf{0.98} & 1.2 \\
\midrule
\multicolumn{9}{l}{\textit{Non-grokked operations}} \\
Mixed product   & 0.03 & 18  & 0.01 & 21   & 0.01 & 16   & 0.04 & 11 \\
Squaring        & 0.00 & 22  & 0.00 & 24   & 0.00 & 19   & 0.00 & 14 \\
\bottomrule
\end{tabular}
\renewcommand{\arraystretch}{1.0}
\caption{Full $2 \times 2$ ablation with reconstruction MSE. Grokked operations
have MSE $\leq 1.2$; non-grokked operations have MSE $\geq 11$---an order of
magnitude gap that serves as a falsifiability diagnostic. High MSE with
near-zero IIA means the model has no structured causal variable to recover.
The plain SAE degrades on harder operations (multiplication 0.72, quartic sum
0.32) because without the structured prior, the sparse encoder has no target
direction. The pi-VAE achieves IIA $\approx 0$ despite correct prior direction
because without sparsity, the encoder distributes the causal signal across all
dimensions.}
\label{tab:ablation_full}
\end{table}

\paragraph{Control experiments.}
Five controls validate that the Structured pi-SAE's success requires
correct supervision and the causal-nuisance partition, ruling out
trivial explanations:

\vspace{0.5em}
\begin{table}[H]
\centering
\small
\begin{tabular}{lcc}
\toprule
\textbf{Control} & \textbf{IIA} & \textbf{Purpose} \\
\midrule
Random-label pi-SAE    & $\leq 0.02$  & Labels carry no causal signal \\
Reconstruction-only ($\alpha{=}0$) & $\leq 0.005$ & No supervision to pin direction \\
Untrained pi-SAE       & 0.000        & Architecture alone is insufficient \\
Plain VAE (no split)   & $\leq 0.08$  & Causal-nuisance partition is necessary \\
Random subspace (DAS)  & $\leq 0.01$  & Baseline for chance-level IIA \\
\bottomrule
\end{tabular}
\caption{Control experiments on grokked multiplication ($p = 113$).
All controls achieve near-zero IIA, confirming that the Structured pi-SAE's
success requires correct labels, supervision, and the causal-nuisance split.}
\label{tab:controls}
\end{table}

\section{Mechanistic analysis: architecture, Fourier dynamics, and spectral compression}
\label{app:mechanistic}

This appendix presents detailed mechanistic analysis of how grokking
creates causal structure. The main text (\S\ref{sec:mechanism})
summarizes these findings and presents the two results most central to
the Grassmannian interpretation (weight-space SVD and subspace
degeneracy).

\subsection{Minimal architecture for grokking}
\label{app:dmf}

What is the simplest model that groks? We compare a depth-2 deep matrix
factorization (DMF) $y = W_2 W_1 x$ against variants with nonlinearity
(DMF + ReLU), a linearized transformer (identity attention + ReLU MLP),
a softmax-only transformer (no MLP), and the full single-layer transformer.
All are trained on modular multiplication ($\Z_{113}$) with weight decay 1.0
for 50{,}000 epochs.

\begin{table}[t]
\centering
\small
\begin{tabular}{lrrl}
\toprule
Architecture & Grok epoch & Fourier align. & Groks? \\
\midrule
DMF + ReLU     & $\sim$2{,}000   & 0.095 & Yes (18$\times$ faster) \\
MLP-only       & $\sim$2{,}500   & 0.99  & Yes \\
Linearized TF  & $\sim$14{,}000  & 0.000 & Yes (no Fourier) \\
Full transformer & $\sim$37{,}000 & 0.10$\to$0.54 & Yes \\
Linear DMF     & ---          & 0.078 & \textbf{No} \\
Softmax-only   & ---          & 0.103 & \textbf{No} \\
\bottomrule
\end{tabular}
\caption{Grokking speed hierarchy. Nonlinearity is necessary (linear DMF never
groks despite nuclear norm dropping from 7{,}340 to 3{,}877). Attention is not
(DMF + ReLU groks 18$\times$ faster). Softmax alone is insufficient.}
\label{tab:grokking_speed}
\end{table}

Three findings emerge. First, \textbf{nonlinearity is necessary}: the linear
DMF never generalizes despite 50{,}000 epochs. Its nuclear norm drops from
7{,}340 to 3{,}877---spectral compression occurs---but test accuracy remains
$\sim$0\%. Compression alone is not sufficient for grokking; the nonlinear
activation function that enables the Fourier representation is essential.

Second, \textbf{attention is not necessary for grokking on modular
multiplication}: DMF + ReLU groks 18$\times$ faster than the full
transformer. The softmax attention mechanism adds enormous overhead to the
grokking process without being required for the underlying algorithm on this
operation. Whether this extends to other operations or multi-task settings
remains open.

Third, the \textbf{linearized transformer groks without Fourier structure}:
it achieves 100\% test accuracy by epoch 14{,}000 but its Fourier alignment
remains at 0.000 throughout. Grokking and Fourier representation are
dissociable---grokking can occur via non-Fourier algorithms. The Fourier
representation is one route to generalization, not the only one.
\citet{xu2026systematic} independently find that the effect of activation
functions on grokking is regime-dependent, consistent with our finding that
nonlinearity is necessary but the specific nonlinearity (ReLU vs.\ softmax
attention) determines which algorithm the model discovers.

\subsection{Fourier structure emerges after generalization}
\label{app:fourier_trajectory}

Prior work identified Fourier representations as the mechanism underlying
grokking in modular arithmetic \citep{nanda2023progress, zhong2024clock}.
A natural assumption is that Fourier structure \emph{drives} grokking:
the model discovers the Fourier basis, then generalizes. Our fine-grained
trajectory analysis (1{,}002 checkpoints at 100-epoch resolution) reveals
the opposite: \textbf{generalization precedes Fourier crystallization}.

\begin{table}[t]
\centering
\small
\begin{tabular}{rrl}
\toprule
Epoch & Test accuracy & Fourier alignment \\
\midrule
37{,}000 & 65.6\% & 0.070 \\
37{,}500 & 91.3\% & 0.093 \\
38{,}000 & 98.6\% & 0.143 \\
38{,}500 & 99.7\% & 0.260 \\
39{,}500 & 99.7\% & 0.498 \\
40{,}000 & 99.8\% & 0.543 \\
\bottomrule
\end{tabular}
\caption{Fine-grained Fourier trajectory for modular multiplication.
Test accuracy crosses 90\% at epoch 37{,}500; Fourier alignment crosses 0.30
approximately 1{,}000 epochs later. The model generalizes \emph{before} its
weights crystallize into the Fourier basis.}
\label{tab:fourier_lag}
\end{table}

Test accuracy crosses 90\% at epoch 37{,}500, but Fourier alignment does not
cross 0.30 until approximately epoch 38{,}500---a lag of $\sim$1{,}000
epochs. The model generalizes before its weights fully organize into the
Fourier basis.

\paragraph{The AGOP double-peak.}
Tracking the Average Gradient Outer Product (AGOP) Fourier alignment through
training reveals an unexpected non-monotonic pattern. For the full transformer:

\begin{table}[t]
\centering
\small
\begin{tabular}{rrrrr}
\toprule
Epoch & Test acc. & AGOP FA & Weight FA & AGOP erank \\
\midrule
5{,}000   & 0.4\%   & 0.108 & 0.080 & 10.0 \\
25{,}000  & 6.3\%   & 0.712 & 0.081 & 9.8 \\
30{,}000  & 19.1\%  & \textbf{0.894} & 0.075 & 8.7 \\
35{,}000  & 95.8\%  & 0.662 (dip) & 0.270 & 5.4 \\
40{,}000  & 100\%   & \textbf{0.974} & 0.864 & 4.7 \\
45{,}000  & 100\%   & 0.888 & 0.977 & 4.4 \\
\bottomrule
\end{tabular}
\caption{AGOP trajectory for modular multiplication. The Fourier alignment of
the gradient outer product shows a double-peak: 0.894 at epoch 30{,}000,
followed by a dip to 0.662 during the grokking transition (epoch 35{,}000),
then crystallization at 0.974 (epoch 40{,}000). The \emph{gradient}
discovers Fourier structure before the \emph{weights} do (weight FA is
0.075 when AGOP FA is 0.894).}
\label{tab:agop}
\end{table}

The AGOP alignment peaks at 0.894, then \emph{dips} to 0.662 during the
actual grokking transition (epoch 35{,}000), before crystallizing at 0.974
post-grokking. The dip coincides with the period of rapid generalization---the
phase transition temporarily disrupts the gradient's Fourier structure before
it stabilizes. By contrast, the DMF's AGOP alignment remains flat at
0.06--0.08 throughout 50{,}000 epochs, confirming that the linear model never
discovers Fourier structure even in its gradients.

The key implication: the gradient ``knows'' the Fourier basis $\sim$10{,}000
epochs before the weights do (AGOP FA is 0.894 at epoch 30{,}000 while weight
FA is still 0.075). Grokking may be the process by which weight-space catches
up to gradient-space structure.

\paragraph{Comparison with memorization-phase encoding.}
\citet{swaroop2026latent} report that latent algorithmic structure is already
encoded during the memorization phase in ReLU MLPs, extractable from
memorization-phase weights. This appears to conflict with our finding that
Fourier alignment lags generalization. The tension is likely
architecture-dependent: in transformers with softmax attention, the Fourier
representation crystallizes in the attention circuit (QK and OV matrices),
which requires coordinated weight reorganization that takes longer than the
MLP-local encoding Swaroop et al.\ observe. The AGOP trajectory supports
this interpretation---the gradient discovers Fourier structure early (epoch
30{,}000), consistent with Swaroop et al., but the weights require an
additional $\sim$10{,}000 epochs to catch up.

\paragraph{Fourier structure lives in attention, not MLPs.}
In a grokked 1-layer model ($p = 113$), we measure Fourier selectivity
across all components: 0 of 512 MLP neurons are Fourier-selective, and the
embedding has a nearly flat Fourier spectrum (participation ratio = 3.93).
The Fourier structure resides entirely in attention: the QK circuit selects
individual frequency pairs (rank-1 structure), while the OV circuit operates
in a broader $k$-dimensional subspace of Fourier modes. This decomposition
is consistent with the ``clock'' interpretation of \citet{zhong2024clock}.

\subsection{Spectral compression during grokking}
\label{app:spectral_compression}

The factored feedforward product matrix $W_\text{FF} = W_\text{in} W_\text{out}$
undergoes dramatic rank collapse during grokking. Tracking the effective rank
(exponential of the Shannon entropy of the normalized singular value spectrum)
through training:

\begin{center}
\small
\begin{tabular}{rrl}
\toprule
Epoch & FF effective rank & Stage \\
\midrule
34{,}000 & 48 & Pre-grokking \\
37{,}500 & 14 & Transition \\
41{,}000 & 9 & Post-grokking \\
\bottomrule
\end{tabular}
\end{center}

The 5.3$\times$ compression from effective rank 48 to 9 confirms that
grokking compresses the model's representational capacity into a
low-dimensional subspace. This compression is necessary but not sufficient:
the linear DMF also achieves spectral compression (nuclear norm
$7{,}340 \to 3{,}877$) without grokking (\S\ref{app:dmf}). What distinguishes
successful grokking is that the compressed subspace acquires
\emph{structure}---specifically, equivariance under the operation's symmetry
group.

\begin{figure}[t]
\centering
\includegraphics[width=\linewidth]{fig3_trajectory}
\caption{Grassmannian trajectory during grokking for modular
multiplication. The Grassmannian distance $d_{\Gr}$ from the current
DAS subspace to the final (fully grokked) subspace decreases as
equivariance rises. The subspace converges to its final position
during the grokking phase transition (dashed line), confirming that
spectral compression and causal structure emerge simultaneously.}
\label{fig:trajectory}
\end{figure}
